\section{Conclusion}\label{sec:conclusion}

Social facts kept in records split along the line the CALM theorem draws through programs. Facts
whose grounding conditions ask only that records exist can be settled by anyone who holds the
records, in any order, and need no authority. Facts that ask for an absence need a seal, a
declaration that a set of records counts as complete, and the seal can be made by the parties the
question concerns as a joint commitment. Coordination is spent once, at the seal, and the fact then
travels as freely as any other.

A seal can be abused like any other power, so the rules that say who may seal have to be records as
well: secondary rules co-signed by the people they bind, frames built in layers by consensus and
customized by groups through additions that stay monotone, changes of any other kind taken through
the legal pathways for adjusting regulation, governing bodies that act through one facade over a
quorum, and appeals that end in institutions outside the collaboration. What remains open is the
content of those rules: who sits on a panel, what quorum may suspend an office, how long a
suspension may take to reach its readers. Those are political choices. The sort shows where they have
to be made, and that a design of records can make them in view of everyone they bind.
